%% Example CV - [Your Name]
%% Structural/formatting reference only - holds no real personal data, so it
%% stays safe to commit and share. Real facts for tailoring a CV live in
%% 01-candidate-profile.md and CLAUDE.md's Candidate Profile section.
%%
%% Compile with: lualatex main_example.tex
%% (pdflatex often fails on modern MiKTeX with fontawesome5 font-expansion errors.)

\documentclass[11pt,a4paper,sans]{moderncv}
\moderncvstyle{banking}
\moderncvcolor{black}

% Force the name and section headings to render in black (color1) - never
% blue. moderncvstylebanking.sty's \colorlet copies (not aliases) the
% pre-scheme accent colour, so the name colours are frozen before
% \moderncvcolor runs; re-let them after. \namefont is the hook every
% name-style macro routes through, so this also works on moderncv 2.3.1
% (Debian/Ubuntu apt), which has no \firstnamestyle/\lastnamestyle at all.
\renewcommand*{\namefont}{\fontsize{34}{36}\bfseries\upshape}
\colorlet{firstnamecolor}{color1}
\colorlet{lastnamecolor}{color1}
\colorlet{namecolor}{color1}
% \MakeUppercase renders every \section{...} heading (Education, Work
% Experience, etc.) in all caps automatically, without retyping each one.
\renewcommand*{\sectionstyle}[1]{{\sectionfont\color{color1}\MakeUppercase{#1}}}

% Times New Roman, matching cover.cls's font in cover_letters/. moderncv's
% "sans" classoption sets \familydefault to \sfdefault at \begin{document},
% so overriding \setsansfont (not just \setmainfont) is what actually reaches
% body text; \setmainfont is set too for any explicit \rmfamily use (e.g. the
% class's bullet-symbol fallback). BoldFont points at the times-bold.ttf
% companion so \bfseries/\textbf (name, section headings, skill/award labels)
% render a true bold face instead of a synthesized/absent-shape substitution.
\usepackage{fontspec}
\setmainfont[Path = OpenFonts/fonts/times-new-roman/, BoldFont = times-bold]{times}
\setsansfont[Path = OpenFonts/fonts/times-new-roman/, BoldFont = times-bold]{times}
% \addressfont drives the contact line under the name (moderncvheadiii.sty
% default: \normalsize\mdseries\upshape). Pinned to an explicit size, matching
% \lettercontent's 11pt/13pt in cover.cls, so the CV and cover letter contact
% lines render identically instead of both happening to use "normal size."
\renewcommand*{\addressfont}{\fontsize{11pt}{13pt}\mdseries\upshape}

\usepackage[utf8]{inputenc}
% pdflatex fallback only (the documented engine is lualatex, which skips this
% branch). Without T1 font encoding pdflatex builds accented letters with
% \accent, and the PDF text layer stores them decomposed - `e` + U+0300 rather
% than U+00E8 - so an ATS keyword match on "Genève" fails while the page looks
% right. moderncv 2.5 loads T1 itself under pdflatex; 2.3.1 (Debian/Ubuntu apt)
% does not. \ifpdftex comes from iftex, which every moderncv version loads.
\ifpdftex\usepackage[T1]{fontenc}\fi
% moderncv loads hyperref itself in an \AtEndPreamble hook, so \hypersetup
% must go in an \AtEndPreamble of our own: on moderncv < 2.4 a top-level
% \usepackage{hyperref} clashes with the class's own
% \RequirePackage[unicode]{hyperref}. From 2.4.0 the class passes its options
% through \PassOptionsToPackage instead, which is what removes that clash.
\AtEndPreamble{\hypersetup{
    colorlinks=true,
    linkcolor=black,
    filecolor=black,
    urlcolor=black,
    pdftitle={[Your Name] - CV},
    % Keep pdfpagemode=UseNone: this block runs after moderncv's own
    % \AtEndPreamble (moderncv.cls sets pdfpagemode there), so a FullScreen
    % value here would win and open every CV in fullscreen presentation mode.
    pdfpagemode=UseNone,
}}
\usepackage[scale=0.80]{geometry}
\usepackage{import}
% Section bullets use plain itemize (not moderncv's own cvitem/cvlistitem),
% which keeps LaTeX's default itemsep/topsep/parsep. Switching to Times New
% Roman's slightly taller line spacing pushed the page from 1 to 2, so
% tighten list spacing globally to keep the whole CV on one page.
\usepackage{enumitem}
\setlist{itemsep=0pt, topsep=2pt, parsep=0pt}

% personal data - bracketed placeholders only; never fill this file in with
% real contact details (see 01-candidate-profile.md / CLAUDE.md instead).
\name{[YOUR}{NAME]}
% No \address/\phone/\email calls: moderncv would print each on the same
% "details" line as \extrainfo anyway (just with icons in front), but the
% cover letter's header has no icons - so the whole contact line is built by
% hand in \extrainfo below, in the exact order/format cover.cls's
% \namesection uses, which is what keeps the two headers pixel-identical.
\extrainfo{[+XX XXXXXXXXXX] | \href{mailto:your.email@example.com}{your.email@example.com} | \href{https://linkedin.com/in/yourprofile}{LinkedIn} | \href{https://yourportfolio.com}{Portfolio} | \href{https://github.com/yourusername}{GitHub}}

\begin{document}

\makecvtitle

% ============================================================
%     PROFILE STATEMENT
% ============================================================

\vspace{6pt}
\small{[Write a 3-5 line profile statement tailored to the role you're applying for. Focus on what makes you uniquely qualified. Example: "Data scientist with a PhD in [field] and X years of industry experience. Combines deep domain expertise in [domain] with strong applied machine learning and Python development skills. Experienced in developing end-to-end ML pipelines, from data ingestion to stakeholder-facing dashboards."]}

% ============================================================
%     EDUCATION
% ============================================================

\section{Education}
\vspace{1pt}
\begin{itemize}

\item{\cventry{[Dates]}{[Degree Name]}{[Institution Name]}{[Location]}{}{\vspace{1pt}
[Concentration, minor, relevant coursework, GPA if strong, honors, etc.]
}}

\end{itemize}

% ============================================================
%     WORK EXPERIENCE
% ============================================================

\section{Work Experience}
\vspace{3pt}
\begin{itemize}

% --- Most Recent Role ---
\item{\cventry{[Dates]}{[Job Title]}{[Company Name]}{[Location]}{}{\vspace{1pt}
\begin{itemize}
    \item {[Responsibility or achievement, with a number where possible]}
\end{itemize}}}

\vspace{3pt}

% --- Previous Role ---
\item{\cventry{[Dates]}{[Job Title]}{[Company Name]}{[Location]}{}{\vspace{1pt}
\begin{itemize}
    \item {[Responsibility or achievement, with a number where possible]}
    \item {[Responsibility or achievement, with a number where possible]}
\end{itemize}}}

\end{itemize}

% ============================================================
%     PROJECTS
% ============================================================

\section{Projects}
\vspace{3pt}
\begin{itemize}

% \cventry bolds its 3rd arg (organization slot), so Project Name goes there
% and Tech Stack goes in the title slot - matching how Company Name (not Job
% Title) renders bold in Work Experience above.
\item{\cventry{[Dates]}{[Tech Stack]}{[Project Name]}{[Link, if any]}{}{\vspace{1pt}
\begin{itemize}
    \item {[What you built and the result/impact, with a number where possible]}
\end{itemize}}}

\end{itemize}

% ============================================================
%     TOOLS & TECHNOLOGIES
% ============================================================

\section{Tools \& Technologies}
\vspace{1pt}
\begin{itemize}

\item \textbf{[Languages]}: [Specific languages. Be concrete.]

\item \textbf{[Frameworks \& Libraries]}: [Specific frameworks, libraries.]

\item \textbf{[Databases]}: [Specific databases, data stores.]

\item \textbf{[Cloud \& DevOps]}: [Platforms, CI/CD, infrastructure tools.]

\item \textbf{[Developer Tools]}: [IDEs, version control, other tooling.]

\end{itemize}

% ============================================================
%     LANGUAGES (optional)
% ============================================================

\section{Languages}
\vspace{1pt}
\begin{itemize}
\item {[Language] ([Proficiency level, e.g. native/fluent]).}
\end{itemize}

% ============================================================
%     PUBLICATIONS (optional - omit this section entirely when there are
%     none, since a placeholder like "[Author(s)]" would fabricate an
%     achievement on a CV meant to stay factual)
% ============================================================

% ============================================================
%     HONORS AND AWARDS (optional)
% ============================================================

\section{Honors and Awards}
\vspace{1pt}
\begin{itemize}
\item{\textbf{[Award or Honor Name]} -- [Brief description, if any] ([Year]).}
\item{\textbf{[Award or Honor Name]}}
\end{itemize}

\end{document}
